% Solutions to the exercises of Chapter 21.
\begin{solutions}
  \item The exponent splits, $e^{-(u^2 + v^2)/2\sigma^2} = e^{-u^2/2\sigma^2}e^{-v^2/2\sigma^2}$, and so does the constant, $\frac1{2\pi\sigma^2} = \frac1{\sqrt{2\pi}\sigma}\cdot\frac1{\sqrt{2\pi}\sigma}$, so
        $G_\sigma(u, v) = g_\sigma(u)g_\sigma(v)$. Sampled on a square $k\times k$ grid the kernel is the outer product $h = gg\T$ of the sampled 1D kernel, and normalizing
        each 1D kernel to unit sum normalizes the product too. Then
        \[
          (h * I)[u, v] = \sum_jg[j]\Big(\sum_ig[i]\,I[u - i, v - j]\Big):
        \]
        the inner sum filters every row with $g$ ($k$ multiplications per pixel) and the outer sum filters the columns of the result ($k$ more), $2k$ in total
        instead of $k^2$. For $\sigma = 1.5$ truncated at $\pm3\sigma$ ($k = 11$) this is 22 instead of 121. Truncating to a \emph{disk} instead of a square would destroy
        the factorization.

  \item The kernel $[a_0, \dots, a_n]$ corresponds to the polynomial $\sum_ia_iz^i$, and convolution of kernels is multiplication of polynomials. Convolving
        $\tfrac12[1, 1]$ with itself $n$ times gives $\big(\tfrac12(1 + z)\big)^n = 2^{-n}\sum_i\binom{n}{i}z^i$: the binomial kernels $\tfrac14[1, 2, 1]$, $\tfrac18[1, 3, 3, 1]$, $\tfrac1{16}[1, 4, 6, 4, 1]$, \dots.
        A non-negative kernel with unit sum is a probability distribution, and the convolution of two distributions is the distribution of the sum of
        independent variables, whose variances add (\cref{ch:prob}). The two-point kernel $\tfrac12[1, 1]$ has variance $\tfrac14$, so the $n$-fold kernel has variance $n/4$:
        $\tfrac12$ for $[1, 2, 1]$, $1$ for $[1, 4, 6, 4, 1]$. By the central limit theorem the binomial kernels approach Gaussians with $\sigma = \sqrt n/2$; OpenCV's
        \code{pyrDown} kernel is a Gaussian with $\sigma = 1$~px.

  \item At pixel $(u, v)$ the correlation sums $S_u[j, i]\,c\,(u + i)$. In each row the weights $-1, 0, 1$ give $c(u + 1) - c(u - 1) = 2c$, and the three rows are
        weighted $1, 2, 1$, so
        \[
          (S_u\star I)[u, v] = (1 + 2 + 1)\cdot2c = 8c,
        \]
        everywhere, for a true derivative $I_u = c$. The factor 8 is the product of the smoothing kernel's sum (4) and the distance between the differenced
        pixels (2). Divide by 8 to obtain the gradient (by 32 for Scharr, whose smoothing kernel sums to 16). With convolution instead of correlation the
        kernel is flipped and the response is $-8c$: check the sign convention of whatever library you use.

  \item A zone plate $I = \tfrac12 + \tfrac12\cos(\pi r^2/n)$ on an $n\times n$ grid (the figure uses $n = 256$) has phase $\pi r^2/n$, hence local frequency $r/n$ cycles per pixel. After
        keeping every fourth pixel the representable limit is $\tfrac18$ cycle per original pixel, so the true rings beyond $r = n/8 = 32$ cannot be represented
        and should become uniform gray; instead, naive sub-sampling folds them into false rings, with new ``centers'' where the frequency is a multiple of the
        new sampling rate $\tfrac14$. A Gaussian blur multiplies frequency $f$ by $e^{-2\pi^2\sigma^2f^2}$. Measuring the false pattern in the region $48 < r < 128$:
        \begin{center}\small
        \begin{tabular}{@{}l c c c c c@{}}
        \toprule
        blur $\sigma$ (px) & 0 & 1 & 1.5 & 2 & 3 \\
        \midrule
        RMS of false pattern & 0.353 & 0.061 & 0.018 & 0.005 & 0.002 \\
        response at $f = \tfrac18$ (kept detail) & 1 & 0.73 & 0.50 & 0.29 & 0.06 \\
        response at $f = \tfrac14$ (aliases to DC) & 1 & 0.29 & 0.06 & 0.007 & $10^{-5}$ \\
        \bottomrule
        \end{tabular}
        \end{center}
        A blur of $\sigma \approx 2$~px, half the shrink factor, removes the aliasing (as in \cref{fig:aliasing-image}); much more also removes detail the smaller image could
        have kept. The same rule, $\sigma \approx s/2$ for a shrink by $s$, gives \code{pyrDown}'s $\sigma = 1$ for $s = 2$.

  \item Along the row $v_0$, linear interpolation gives $(1 - a)I_{00} + aI_{10}$, and along the row $v_0 + 1$, $(1 - a)I_{01} + aI_{11}$. Interpolating these linearly in $v$ with
        weights $1 - b$ and $b$ gives the formula of \cref{sec:interpolation}. In the other order, the columns give $(1 - b)I_{00} + bI_{01}$ and $(1 - b)I_{10} + bI_{11}$, and
        interpolating in $u$ gives $(1 - a)(1 - b)I_{00} + (1 - a)bI_{01} + a(1 - b)I_{10} + abI_{11}$, the same four products. Without expanding: both results have the form
        $c_0 + c_1a + c_2b + c_3ab$ and take the four corner values at the corners, and those four conditions determine the four coefficients uniquely. The
        interpolant is linear along every row and column but quadratic along diagonals, and its derivatives jump at pixel boundaries.

  \item Fit $r(x) = Ax^2 + Bx + C$ through $(-1, r_{-1})$, $(0, r_0)$, $(1, r_1)$: $C = r_0$, $A + B + C = r_1$ and $A - B + C = r_{-1}$, so
        $B = \tfrac12(r_1 - r_{-1})$ and $A = \tfrac12(r_{-1} - 2r_0 + r_1)$. The vertex is at
        \[
          \delta = -\frac{B}{2A} = \frac{r_{-1} - r_1}{2(r_{-1} - 2r_0 + r_1)} .
        \]
        If $r_0$ is the largest of the three, $A < 0$ (a maximum), and $|r_{-1} - r_1| \le (r_0 - r_{-1}) + (r_0 - r_1) = 2|A|$, so $|\delta| \le \tfrac12$: the refined peak stays within
        the pixel. Real peaks are not parabolas, and the fit is biased toward the pixel center: for a Gaussian peak of width $\sigma = 1$~px at a true
        offset of $0.3$~px it returns $0.25$, a bias of up to $0.05$~px, larger for sharper peaks.
\end{solutions}
